% In this file you should put macros to be used only by
% the printed version. Of course they should have a corresponding
% version in macros/web.tex.
% Typically the printed version could have more fancy decorations.
% This should be a very short file.
%
% This file starts with dummy macros that ensure the pdf
% compiler will ignore macros provided by plasTeX that make
% sense only for the web version, such as dependency graph
% macros.


% Dummy macros that make sense only for web version.
\newcommand{\lean}[1]{}
\newcommand{\discussion}[1]{}
\newcommand{\leanok}{}
\newcommand{\mathlibok}{}
\newcommand{\notready}{}
% Make sure that arguments of \uses and \proves are real labels, by using invisible refs:
% latex prints a warning if the label is not defined, but nothing is shown in the pdf file.
% It uses LaTeX3 programming, this is why we use the expl3 package.
\ExplSyntaxOn
\NewDocumentCommand{\uses}{m}
 {\clist_map_inline:nn{#1}{\vphantom{\ref{##1}}}%
  \ignorespaces}
\NewDocumentCommand{\proves}{m}
 {\clist_map_inline:nn{#1}{\vphantom{\ref{##1}}}%
  \ignorespaces}
\ExplSyntaxOff

% Fallbacks for math delimiters missing from the print font setup.
\providecommand{\llbracket}{[\mkern-4mu[}
\providecommand{\rrbracket}{]\mkern-4mu]}

% Long Lean names in typewriter font: allow a line break after each dot and each underscore
% (a name is never hyphenated), and let TeX stretch the spaces of a paragraph further rather
% than push a long name or formula into the margin.
\let\blueprintorigtexttt\texttt
\ExplSyntaxOn
\tl_new:N \l__blueprint_tt_tl
\RenewDocumentCommand{\texttt}{m}
 {
  \tl_set:Nn \l__blueprint_tt_tl {#1}
  \tl_replace_all:Nnn \l__blueprint_tt_tl {.} {.\allowbreak}
  \tl_replace_all:Nnn \l__blueprint_tt_tl {\_} {\_\allowbreak}
  \blueprintorigtexttt{\l__blueprint_tt_tl}
 }
\ExplSyntaxOff
\setlength{\emergencystretch}{3em}
